\scanpage{5}
\[
\operatorname{réal}(S,A,R)\lhook\joinrel\longrightarrow E
\]
telle que l'image d'un arc soit un segment.

\uline{Prop} Pour qu'un graphe fini admette une réalisation géom. rectiligne, (dans un vectoriel convenable sur $\mathbb{R}$) il f et s qu'il \deleted{a} soit strict -- et alors il admet une réalisation rectiligne dans $\mathbb{R}^{3}$. De façon précise, si $S\hookrightarrow\deleted{\unsure{\mathbb{R}^{3}|R^{3}}}E$ \inserted{$(\simeq\mathbb{R}^{3},\ \unsure{\text{v. ex.}|\text{p. ex.}|\text{v. a.}})$} est une immersion telle que
% Pencil drawing in the left margin; the small lower stroke is detached.
\margin{\figure{\begin{tikzpicture}[x=.025cm,y=.025cm,line cap=round,line join=round]
\draw[line width=.65pt] (0,0) ellipse (1.8 and 1.3);
\draw[line width=.65pt] (-.1,-2.2)--(-.4,-3.6);
\draw[line width=.75pt] (1.4,1.1)--(31.7,34.4);
\fill (31.1,34.1)--(30.7,37.3)--(33.8,35.9)--(33.3,33.8)--cycle;
\draw[line width=.35pt] (35.4,39)--(59.8,65.3);
\fill (59,65.5)--(60.1,68)--(60.7,65)--cycle;
\end{tikzpicture}}}
\[
\forall x,y\in S\quad x\ne y,\quad \text{on ait}\quad
\deleted{\unsure{\{|\exists}}
\underbrace{\overline{x,y}}_{\substack{\text{segm}\\\text{joignant}\\x\ \text{à}\ y}}\cap S=\{x,y\},
\]
\inserted{(i.e. trois points \inserted{distincts} jamais alignés)}
\inserted{et $\forall x,y,z,t$ distincts, $(\overline{x,y})\cap(\overline{z,t})=\varnothing$}

$\exists 1$ réalis géom. rectiligne $\unsure{\operatorname{réll}|\operatorname{réel}}(S,A,R)\lhook\joinrel-E$ qui $f$ \unsure{prolonge|prolonge\textsuperscript{?}}\ldots

Graphe connexe $=$ réal. géom \unsure{connex|connexe} $=$ qui n'est pas \deleted{\unsure{réunion}} ``somme'' de deux autres graphes $\ne\varnothing$

\uline{NB} Tout graphe est de façon unique réunion de graphes connexes \inserted{et disjoints}

[ graphe somme : $S\amalg S',\quad A\amalg A',\quad R\amalg R'\ \ldots$ ]

% A curved arrow in the left margin leads from NB to Graphe polygonal.
\margin{\figure{\begin{tikzpicture}[baseline=-.5ex,x=.75cm,y=.75cm]\draw[->] (.70,1.8) .. controls (-.30,.60) and (0,.25) .. (1.8,0);\end{tikzpicture}}}
% multi is circled above polygonal; the contour qualifier stands above the bracketed formula.
% The heading and the bracketed formula share a baseline.
\[
\text{Graphe }\overset{\inserted{\text{multi}}}{\text{polygonal}}
\qquad
\overset{\inserted{\text{(ou \uline{contour} (combinatoire))}}}{\bigl[\quad\forall s\in S,\quad\operatorname{card}R\{s\}=2\quad\bigr]}
\]

Dualité des \deleted{graphes polygonaux} contours (combinatoires).

\deleted{\unsure{donq|donc}} les composantes \inserted{\unsure{conn|conn.|connexes}} d'un \deleted{graphe polygonal}\inserted{contour} sont \deleted{polygonales.} des contours.

Classification des \deleted{graphes polygonaux}\inserted{contours} conn. $\ne\varnothing$

\uline{Prop} Un \deleted{graphe polyg.}\inserted{contour} connexe \inserted{$\ne\varnothing$} est connu à isom. près quand on connait $\operatorname{card}S$ ($=\operatorname{card}A$), qui est n'importe quel cardinal \deleted{fini ou infini}
\[
2\le c\le\aleph_{0}
\]
(strict \unsure{sss|ssi} $c\ne2$ i.e. $c\geqslant3$)
% archive sheet number: 4
